\documentclass[a4paper, fleqn]{article}
\usepackage{enumitem}
\usepackage{amsmath, amssymb}
\usepackage{xstring}
\usepackage{graphicx} % For including figures
\graphicspath{{./figs/}} % Path to figures
\usepackage[labelformat=empty]{caption} % To remove figure labels
% Set engineering notation
\providecommand{\sci}[1]{\protect\ensuremath{\times 10^{\StrSubstitute[0]{#1}{e}{}}}}
\setlength{\parindent}{0pt} % Set paragraph indentation to zero

\begin{document}

\underline{\textbf{Tutorial 5: Belts}}
\vspace{10pt}

\section*{Theory}
\begin{enumerate}
    \item What are the common types of belts used in mechanical systems?
    \item What are the advantages and disadvantages of using belts for power transmission?
    \item How do you determine the appropriate belt size for a specific application?

\end{enumerate}


\section*{Calculations}

\textbf{Question 1}

A plastic flat belt 60 mm wide and 0.5 mm thick transmits 10 kW. The input pulley has a diameter of 300 mm, it rotates at 2800 rpm, and the output pulley speed is 1600 rpm; the pulleys are 700 mm apart, the coefficient of friction is 0.2, and belt weight is 25 kN/m3. It is assumed that the driver is a high torque motor, and the driven machine is under a medium shock load. Determine;

\begin{enumerate}[label=(\roman*)]
    \item the torque at smaller pulley
    \item the contact angle at driver and driven pulleys
    \item the centrifugal tension in the belt
    \item the maximum tension and stress in the belt
\end{enumerate}
 
\vspace{10pt}
\textbf{Example Solution}
\vspace{10pt}

\textbf{Given Data:}


Belt width, b = 60 mm = 0.06 m

Belt thickness, t = 0.5 mm = 0.0005 m

Power, P = 10 kW = 10000 W

Diameter of smaller pulley, d1 = 300 mm = 0.3 m

Speed of smaller pulley, n1 = 2800 rpm

Speed of larger pulley, n2 = 1600 rpm

Distance between pulleys, c = 700 mm = 0.7 m

Coefficient of friction, f = 0.2

Belt weight, = 25 kN/m

\vspace{10pt}
i - Torque at smaller pulley

\begin{equation*}
    \begin{aligned}
    Power, P &= \dfrac{2\pi n_{rpm}}{60} T\\
    T &= \dfrac{(10000) 60}{2 \pi (2800)}\\
    &=34.1 Nm
    \end{aligned}
\end{equation*}

ii - Contact angle

\begin{equation*}
    \begin{aligned}
    \sin \alpha &= \dfrac{r_2 - r_1}{c}\\
    \end{aligned}
\end{equation*}

Diameter of larger pulley is unknown, thus we can find it using the relation between the speeds of the pulleys.
\begin{equation*}
    \begin{aligned}
    \dfrac{n_1}{n_2} &= \dfrac{d_2}{d_1}\\
    d_2 &= \dfrac{n_1 d_1}{n_2} = \dfrac{2800 \times 0.3}{1600} = 0.525 m
    \end{aligned}
\end{equation*}

where, smaller pulley radius,$r_1 = \dfrac{0.3}{2} = 0.15 m$
\vspace{10pt}

larger pulley radius, $r_2  = \dfrac{0.525}{2} = 0.2625 m$\\

Find $\alpha$:
\begin{equation*}
    \begin{aligned}
    \sin \alpha &= \dfrac{0.2625 - 0.15}{0.7}\\
    \alpha &= 9.25^{\circ}\\
    \end{aligned}
\end{equation*}

Equation to find contact angle is in radians $\phi = \pi - 2\alpha$, thus we need to
convert to radians:
\begin{equation*}
    \begin{aligned}
    \alpha &= 9.25^{\circ} \times \dfrac{\pi}{180} = 0.1614 rad
    \end{aligned}
\end{equation*}

Contact angle for driver pulley(smaller pulley):
\begin{equation*}
    \begin{aligned}
    \theta_1 &= \pi - 2\alpha\\
    &= \pi - 2(0.1614) = 2.82 rad
    \end{aligned}
\end{equation*}

Contact angle for driven pulley(larger pulley):
\begin{equation*}
    \begin{aligned}
    \theta_2 &= \pi + 2\alpha\\
    &= \pi + 2(0.1614) = 3.48 rad
    \end{aligned}
\end{equation*}

iii - Centrifugal tension in the belt

Centrifugal tension, $F_c = \dfrac{w}{g} v^2$, or {$F_c = \rho A v^2$}

where $w$ is belt weight per unit length (N/m) and $v$ is velocity of belt. $g$ is acceleration due to gravity.

Belt weight given as $25 kN/m^3 = 25000 N/m^3$. Need to find belt weight per unit length, $w$ in (N/m): 

\begin{equation*}
    \begin{aligned}
    w &= 25000 \left(\dfrac{N}{m^3}\right) \times \text{Cross section area of belt} (m^2) \\
    &= 25000 \left(\dfrac{N}{m^3}\right) \times 0.06m \times 0.0005m = 0.75 N/m
    \end{aligned}
\end{equation*}

where,

Belt width, b = 60 mm = 0.06 m

Belt thickness, t = 0.5 mm = 0.0005 m

Next, velocity of belt, $v = \dfrac{\pi d_1 n_1}{60} = \dfrac{\pi \times 0.3 \times 2800}{60} = 43.98 m/s$\\

Now, calculate centrifugal tension, $F_c$:

\begin{equation*}
    \begin{aligned}
    F_c &= \dfrac{0.75}{9.81} (43.98)^2\\
    &= 147.88 N
    \end{aligned}
\end{equation*}

iv - Maximum tension and stress in the belt

Maximum tension is tension of the belt at tight side, $F_1$

Using belt tension equation for flat belt:
\begin{equation*}
    \begin{aligned}
    \dfrac{F_1-F_c}{F_2-F_c} &= e^{f \phi}\\
    \end{aligned}
\end{equation*}

where, 

$F_1$ : tension on the tight side in the belt, 

$F_2$ : the tension on the slack side of the belt,

$F_c$ : centrifugal tension, 

$f$ : coefficient of friction, 

$\phi$ : contact angle in radians.

\begin{equation*}
    \begin{aligned}
    \dfrac{F_1-147.88}{F_2-147.88} &= e^{0.2(2.8188)}\\
    \dfrac{F_1-147.88}{F_2-147.88} &= 1.7573\\
    F_1 - 147.88 &= 1.7573 (F_2 - 147.88)\\
    F_1 &= 1.7573 F_2 - 259.89 + 147.88\\
    \end{aligned}
\end{equation*}

\begin{equation}
    \begin{aligned}
    F_1 &= 1.7573 F_2 - 112.01\\
    \end{aligned}
    \label{eq1}
\end{equation}

From torque equation:
\begin{equation*}
    \begin{aligned}
    T &= (F_1 - F_2) r_1\\
    34.1 &= (F_1 - F_2) (0.15)\\
    \dfrac{34.1}{0.15}  &= F_1 - F_2\\
    227.33  &= F_1 - F_2\\
    \end{aligned}
\end{equation*}

\begin{equation}
    \begin{aligned}
    F_2 &= F_1 - 227.33\\
    \end{aligned}
    \label{eq2}
\end{equation}

Substitute equation 2 into equation 1:
\begin{equation*}
    \begin{aligned}
    F_1 &= 1.7573 (F_1 - 227.33) - 112.00\\
    F_1 &= 1.7573 F_1 - 399.5 - 112.00\\
    F_1 - 1.7573 F_1 &= -511.5\\
    -0.7573 F_1 &= -511.5\\
    F_1 &= \dfrac{511.5}{0.7573}\\
    &= 675.42 N
    \end{aligned}
\end{equation*}

OR use belt tension relationship:

\begin{equation*}
    \begin{aligned}
    F_1 &= Fc+ \left(\dfrac{\gamma}{\gamma-1}\right) \dfrac{T_1}{r_1}\\
    \end{aligned}
\end{equation*}

where,

$\gamma = e^{f \phi}$ for flat belt.

Substituting the values:
\begin{equation*}
    \begin{aligned}
        \gamma &= e^{0.2(2.8188)} = 1.7573\\
        F_1 &= 147.88 + \left(\dfrac{1.7573}{1.7573-1}\right) \dfrac{34.1}{0.15}\\
        F_1 &= 675.48 N
    \end{aligned}
\end{equation*}

The final answer is approximately the same using both methods.

\vspace{10pt}
Maximum stress in the belt, A is the cross-sectional area of the belt:
\begin{equation*}
    \begin{aligned}
    \sigma_{max} &= \dfrac{F_1}{A}\\
    \sigma_{max} &= \dfrac{675.48}{0.06 \times 0.0005}\\
    &= 22.52 MPa
    \end{aligned}
\end{equation*}

\newpage
\textbf{Question 2}

V belt drive has a 200mm diameter small sheave with 170° contact angle, 38° included angle, coefficient of friction of 0.15, 1600rpm driver speed, belt weight 8N/m and a tight side tension of 3kN. Determine the power capacity of the drive.

\newpage

\textbf{Question 3}

A v belt drives has an included angle of $38^\circ$, belt density, $\rho=2100 \, \text{kg/m}^3$, belt cross sectional are of $145 mm^2$, coefficient of friction of 0.25, $r_1$=150mm, $\phi$=160, $n_1$=3000rpm and $F_2$=800N. Calculate

\begin{enumerate}[label=(\roman*)]
    \item the maximum power transmitted
    \item the maximum stress in the belt
\end{enumerate}

\newpage

\textbf{Question 4}

In a horizontal V belt drive for a centrifugal blower, the blower is belt driven at 600 rpm by a 15 kW, 1750rpm electric motor. The groove angle of the pulley is 35°. The center distance is twice the diameter of the larger pulley and velocity of the belt is 20 m/s. The density of the belt material is 1500kg/m3 and maximum allowable stress is 4MPa. The coefficient of friction of motor pulley is 0.5. Determine

\begin{enumerate}[label=(\roman*)]
    \item the diameter of the blower pulley
    \item the belt length
    \item the size of the belt in mm2
\end{enumerate}

\newpage

\textbf{Question 5}

A motor with 16 kW power is attached to a compressor via a flat belt with a length of 3 m.
The belt experienced a maximum stress of 4 MPa in an operation and medium shock with
Ks =1.5. If the coefficient of friction between the belt surface and the pulleys are 0.4 and the density of the belt is 950 kg/m3 and the details of the belt drive are tabulated in Table 1. Determine;

\begin{table*}[h]
\centering
\begin{tabular}{|c|c|c|}
\hline
 & Diameter & Speed \\
\hline
 Compressor Pulley & 0.6 & 1200 \\
 Motor Pulley & dm & 1800 \\
\hline
\end{tabular}
\caption{Table Q5}
\end{table*}

\begin{enumerate}[label=(\roman*)]
    \item the diameter of the motor pulley, dm
    \item the center distance, c between the pulleys.
    \item angle of wrap or contact angle, of both pulleys. Does the angle equal? In what circumstances the angles will be equal?
    \item torque on both pulleys. Does the torque on both pulleys equal? Why? In what circumstances the torque will be equal? How about power?
    \item Velocity of the belt, v
    \item Centrifugal force, Fc
    \item Cross section of belt, A
    \item Tight tension on belt, F1
    \item Slack tension on belt, F2
    \item The initial tension on the belt, Fi
\end{enumerate}


\newpage

\textbf{Answer}

\textbf{Q2}
\begin{enumerate}[label=(\roman*)]
    \item $P = 34.6 kW$
\end{enumerate}

\textbf{Q3}
\begin{enumerate}[label=(\roman*)]
    \item $P = 44.0 kW$
    \item $\sigma_{max} = 11.96 MPa$
\end{enumerate}

\textbf{Q4}
\begin{enumerate}[label=(\roman*)]
    \item $d_2 = 0.636 m = 636 mm$
    \item $L = 3.92 m$
    \item $A = 223 mm^2$
\end{enumerate}

\textbf{Q5}
\begin{enumerate}[label=(\roman*)]
    \item $d_1 = 0.4 m = 400 mm$
    \item $c = 0.708 m$
    \item $\phi_1 = 163.76^{\circ}$, $\phi_2 = 196.24^{\circ}$
    \item $T_1 = 84.88 Nm$, $T_2 = 127.32 Nm$
    \item $v = 37.70 m/s$
    \item $F_c = 1350225.5 A$, Fc = 638.66 N (After find A)
    \item $A = 473 mm^2$
    \item $F_1 = 1261.42 N$
    \item $F_2 = 837.19 N$
    \item $F_i = 1049.31 N$
\end{enumerate}


\end{document}